% from ega4-20-fr.tex
\oldpage[IV-4]{235}
et montrer qu'il est injectif (resp. bijectif sous les conditions de~(ii)).
Or, on peut écrire $S^{-1}A=\varinjlim A_t$, où $t$ parcourt l'ensemble des éléments \emph{réguliers} de $A$, ordonné suivant la relation « $t$ est un diviseur de $t'$ » \sref[0]{0.1.4.5} ; on peut aussi écrire $A_t=\Gamma(\mathrm{D}(t),\sh{O}_X)$.
D'autre part on a par définition $M'(X)=\varinjlim\Gamma(U,\sh{O}_X)$, où $U$ parcourt l'ensemble des ouverts schématiquement denses de $X$ (ordonné par la relation $\supset$), et l'application \sref{IV.20.2.11.1} n'est autre que l'application canonique provenant du fait que les $\mathrm{D}(t)$ constituent une partie de l'ensemble des ouverts schématiquement denses dans $X$ \sref{IV.20.2.9}.
Notons que par définition d'un ouvert schématiquement dense, l'homomorphisme $\Gamma(U,\sh{O}_X)\to\Gamma(U',\sh{O}_X)$ pour deux tels ouverts $U'\subset U$ est toujours \emph{injectif}, donc il en est de même des homomorphismes $\Gamma(U,\sh{O}_X)\to M'(X)$, et on en conclut que \sref{IV.20.2.11.1} est injectif.
Pour prouver que \sref{IV.20.2.11.1} est bijectif, il suffit de montrer que l'ensemble des $\mathrm{D}(t)$ est \emph{cofinal} à l'ensemble des ouverts schématiquement denses, autrement dit que pour un tel ouvert $U$, il existe $t$ régulier dans $A$ tel que $\mathrm{D}(t)\subset U$.
Or, lorsque $X$ est réduit et que les composantes irréductibles de $X$ forment un ensemble localement fini, cet ensemble est fini puisque $X$ a été supposé affine, autrement dit $A$ n'a qu'un nombre fini d'idéaux premiers minimaux $\mathfrak{p}_i$ ; comme $A$ est réduit, l'intersection des $\mathfrak{p}_i$ est réduite à $0$, et dire que $t$ est régulier équivaut donc à dire que $t$ n'appartient à aucun des $\mathfrak{p}_i$ ; on conclut donc par le raisonnement de \sref[I]{I.7.1.9.1}.
Lorsque $A$ est noethérien, dire que $U=X-Y$ (où $Y=V(\mathfrak{J})$ est fermé dans $X$) est schématiquement dense signifie \sref{IV.5.10.2} que $Y$ \emph{ne rencontre pas} $\operatorname{Ass}(\sh{O}_X)$, et en vertu de Bourbaki, \emph{Alg.~comm.}, chap.~IV, \textsection1, n\textsuperscript{o}~4, prop.~8, cela entraîne l'existence d'un $t\in\mathfrak{J}$ tel que $t$ soit $A$-régulier, donc $U\supset\mathrm{D}(t)$.

On a en outre prouvé au cours de cette démonstration le
\end{proof}

\begin{corollary}[20.2.12]
\label{IV.20.2.12}
Si $X=\Spec(A)$, où $A$ est noethérien, ou réduit et n'ayant qu'un nombre fini d'idéaux premiers minimaux, les ouverts schématiquement denses dans $X$ sont ceux qui contiennent un ouvert de la forme $\mathrm{D}(t)$, où $t$ est régulier dans $A$, et $M(X)=M'(X)$ est l'anneau total des fractions $S^{-1}A$, où $S$ est l'ensemble des éléments réguliers de $X$.
\end{corollary}

\begin{remarks}[20.2.13]
\EGAReaderTarget{ega4.c04.statement.20-2-13}{IV.20.2.13}
\begin{enumerate}[label=(\roman*)]
  \item Soient $\vphi$ un élément de $M(X)$, $\vphi'$ son image dans $M'(X)$ ; on a évidemment, par définition (\sref{IV.20.1.4} et \sref{IV.20.2.8}) $\operatorname{dom}(\vphi)\subset\operatorname{dom}(\vphi')$ ; mais en fait on a l'égalité $\operatorname{dom}(\vphi)=\operatorname{dom}(\vphi')$, car il existe une section de $\sh{O}_X$ au-dessus de $\operatorname{dom}(\vphi')$, appartenant à la classe $\vphi'$, et la fonction méromorphe correspondante est égale à $\vphi$ \sref{IV.20.2.11}[(i)], donc $\operatorname{dom}(\vphi')\subset\operatorname{dom}(\vphi)$.

  \item On a déjà noté que lorsque $X$ est \emph{réduit}, on a $\sh{M}'_X=\sh{R}(X)$ (faisceau des fonctions rationnelles sur $X$ \sref[I]{I.7.3.2}) ; si en outre les composantes irréductibles de $X$ forment un ensemble localement fini, on a $\sh{M}_X=\sh{M}'_X=\sh{R}(X)$.
  En général, puisque tout ensemble ouvert schématiquement dense est dense, on a un homomorphisme canonique $\sh{M}'_X\to\sh{R}(X)$, mais même lorsque $X$ est localement noethérien, cet homomorphisme n'est pas nécessairement injectif.
  Par exemple, si $X=\Spec(A)$, où $A$ est un anneau noethérien tel que $\operatorname{Ass}(A)$ contienne des idéaux premiers immergés (ce qui entraîne que $A$ n'est pas réduit), on a vu que $M(X)$ et $M'(X)$ s'identifient à l'anneau total de fractions $S^{-1}A$,
  